\chapter{Conclusiones}

\section{Estado de las correcciones}

Las siete tareas de auditoría mBERT (T-01 a T-07) están \textbf{completadas} con evidencia reproducible:

\begin{enumerate}[noitemsep]
  \item Baseline mayoritaria corregida (0.128).
  \item Checkpoint empaquetado y verificable (F1 0.7777).
  \item McNemar pareado ejecutado (equivalencia estadística vs XLM-R).
  \item Métrica oficial única con trazabilidad completa.
  \item Cortes y plataforma consolidados desde predicciones oficiales.
  \item Robustez adversarial v1 medida; mitigación B2 rechazada por tradeoff en test.
  \item Batería extendida 144 casos (A--G) integrada; comparación con referencia XLM-R.
\end{enumerate}

\section{Riesgos residuales}

\begin{itemize}[noitemsep]
  \item \textbf{Amenazas ofuscadas:} 6/11 correctas en batería adversarial; FN elevados en amenazas.
  \item \textbf{Emojis:} 7/18 correctos; asociación espuria emoji--toxicidad persiste.
  \item \textbf{Sesgo dialectal:} FPR alto en jerga salvadoreña inofensiva (familia D y sonda en test).
  \item \textbf{Plataforma Facebook:} F1-macro inferior a X; requiere más datos o ajuste de umbral por slice.
\end{itemize}

\section{Recomendaciones}

\begin{enumerate}[noitemsep]
  \item Ampliar corpus con amenazas ofuscadas y textos no tóxicos con jerga/emojis.
  \item Mantener evaluación adversarial versionada en cada iteración del modelo.
  \item No adoptar mitigaciones que degraden el F1 oficial sin análisis de costo-beneficio por slice.
  \item Presentar mBERT y XLM-R como estadísticamente equivalentes en test ($n=460$).
\end{enumerate}

\section{Entrega a Cristian}

Este paquete (\texttt{correcciones-export/}) contiene todo lo necesario para auditoría, defensa de tesis y replicación en Colab, sin incluir el trabajo XLMR ni las correcciones de Julio.
